% Part: normal-modal-logic
% Chapter: syntax-and-semantics
% Section: introduction

\documentclass[../../../include/open-logic-section]{subfiles}

\begin{document}

\olfileid{nml}{syn}{int}

\olsection{Pambuka}

Logika modal ngrembug \emph{proposisi modal} lan relasi konsekuensi
ing antarane. Tuladha proposisi modal yaiku:
\begin{enumerate}
\item Mesthi yen $2+2=4$.
\item Mesthi bisa kelakon yen sesuk udan.
\item Yen mesthi bisa kelakon~$!A$, mula bisa kelakon~$!A$.
\end{enumerate}
Kemungkinan lan keniscayan dudu siji-sijine modalitas: konektif
unari liyane uga digolongake dadi modalitas, upamane, ``kudune~$!A$,''
``Bakal kelakon yen~$!A$,'' ``Dana ngerti yen~$!A$,''
utawa ``Dana percaya yen~$!A$.''

Logika modal sepisanan katon ing karya Aristoteles \emph{De
  Interpretatione}: dheweke kang sepisanan nggatekake yen keniscayan
ngimplikasi kemungkinan, nanging ora kosok baline; yen kemungkinan
lan keniscayan bisa padha-padha ditemtokake saka siji lan sijine;
yen $!A \land !B$ bisa bener, mula $!A$ bisa bener lan $!B$
bisa bener, nanging ora kosok baline; lan yen $!A \to !B$
mesthi bener, banjur yen $!A$ mesthi bener, $!B$ uga mesthi bener.

Pendekatan modern kapisan marang logika modal yaiku karya
C.~I. Lewis, kang pungkasane ngasilake karya Lewis lan Langford,
\emph{Symbolic Logic} (1932). Lewis \& Langford ora marem karo
representasi implikasi nganggo kondisional material: $!A \lif !B$
dudu sulih kang becik kanggo ``$!A$ ngimplikasi $!B$.'' Mula,
padha ngusulake supaya implikasi ditegesi minangka ``Mesthi, yen $!A$
banjur $!B$,'' dilambangake dadi $!A \strictif !B$. Nalika nyoba
njlentrehake sifat-sifat kang beda, Lewis nemtokake limang sistem
modal, \Log{S1}, \ldots, \Log{S4}, \Log{S5}; loro kang pungkasan
isih dienggo.

Pendekatan Lewis lan Langford murni \emph{sintaktis}: padha
nemtokake aksioma lan aturan kang lumrah, banjur nliti apa kang bisa
dibuktekake kanthi sarana iku. Pendekatan semantik suwe ora ketemu,
nganti Rudolf Carnap ngupaya sepisanan ing \emph{Meaning and Necessity}
(1947) kanthi nggunakake gagasan \emph{deskripsi kahanan}, yaiku
kumpulan ukara atomik (kang ``bener'' ing deskripsi kahanan kasebut).
Sawise definisi kabeneran ditrapake marang ukara sembarang $!A$,
Carnap netepake $!A$ minangka \emph{mesthi bener} yen bener ing
kabeh deskripsi kahanan. Pendekatan Carnap ora bisa nangani modalitas
\emph{bertingkat}: ukara arupa ``Bisa wae mesthi \ldots bisa wae
$!A$'' mesthi nyusut dadi modalitas paling njero.

Terobosan utama ing semantik modal teka saka artikel Saul
Kripke ``A Completeness Theorem in Modal Logic'' (JSL 1959). Kripke
ngedhasarake karyane marang gagasan Leibniz yen sawijining pratelan
mesthi bener yen bener ``ing kabeh donya kang mungkin.'' Nanging,
gagasan iki duwe kekurangan kaya pendekatan Carnap: kabeneran
pratelan ing sawijining donya $w$ (utawa deskripsi kahanan $s$)
ora gumantung marang $w$ babar pisan. Mula Kripke nganggep yen
donya-donya dihubungake dening \emph{relasi aksesibilitas} $R$,
lan yen pratelan arupa ``Mesthi $!A$'' bener ing donya $w$ yen
lan mung yen $!A$ bener ing kabeh donya $w'$ kang \emph{bisa
diakses saka} $w$. Semantik kang migunakake sawijining wujud
pendekatan iki diarani semantik Kripke lan ngidini logika modal
(ing wangun jamak) berkembang kanthi cepet.

Nalika ditafsirake nganggo semantik Kripke, logika modal
nuduhake kepriye \emph{struktur relasional} katon ``saka njero.''
Struktur relasional iku himpunan kang dilengkapi relasi biner
(umpamane, himpunan siswa ing sawijining kelas kang diurutake miturut
nomer jaminan sosiale yaiku struktur relasional). Nanging,
struktur relasional sejatine ana ing maneka warna domain: saliyane
kemungkinan relatif antarane kahanan donya, ana kahanan epistemik
sawijining agen kang gegandhengan miturut kemungkinan epistemik,
utawa kahanan sistem dinamis kanthi transisi kahanane, lan
sapanunggalane. Logika modal bisa kanggo ngemodelake kabeh iki:
sing kapisan ngasilake logika modal biasa, yaiku logika aletik;
liyane ngasilake logika epistemik, logika dinamis, lan sapanunggalane.

Kita nggatekake siji sudut pandang tartamtu, kang dikenal
dening ahli logika modal minangka ``teori korespondensi.'' Salah siji
temuan awal Kripke kang wigati yaiku akeh sifat relasi
aksesibilitas~$R$ (kayata transitif, simetris, lan sapanunggalane)
bisa dicireni \emph{ing basa modal} dhewe nganggo ``skema modal''
kang cocog. Upamane, ahli logika modal kandha yen refleksivitas
$R$ ``korespondensi'' karo skema ``Yen mesthi~$!A$, banjur~$!A$''.
Utamane kita nliti teori korespondensi sawetara sistem klasik
logika modal (umpamane, \Log{S4} lan \Log{S5}) kang diasilake saka
kombinasi skema \Ax{D}, \Ax{T}, \Ax{B}, \Ax{4}, lan~\Ax{5}.

\end{document}
